\chapter{Faecal Incontinence}
\index{Faecal Incontinence}

\section*{Definition and Overview}
Faecal incontinence is the involuntary leakage of stool or wind from the bowel, and it is sometimes called bowel incontinence. It ranges from occasional slight leakage to the complete inability to control bowel movements. Faecal incontinence is common, particularly in older people and after childbirth, but it is greatly under-reported because of embarrassment. It has a major impact on quality of life, yet it is highly treatable, and no one should accept it without seeking help.

\section*{Epidemiology}
Faecal incontinence affects a significant proportion of adults, with the prevalence increasing with age and in residential care. It is more common in women, partly because of childbirth-related pelvic floor damage, and it is common in people with chronic bowel conditions, neurological disease, and dementia. Many people do not tell their doctor, so the true rate is higher than reported, and the condition is a common cause of social withdrawal and admission to care.

\section*{Aetiology and Pathophysiology}
Bowel control depends on muscles, nerves, and the rectum working together.
\begin{itemize}
    \item \textbf{Anal sphincter weakness:} damage to the muscles around the anus, often from childbirth, surgery, or ageing
    \item \textbf{Nerve damage:} injury to the nerves controlling the bowel, from childbirth, spinal problems, diabetes, or surgery
    \item \textbf{Diarrhoea:} loose stool overwhelms the normal control mechanisms, causing leakage
    \item \textbf{Constipation and impaction:} hard stool can cause overflow leakage of liquid stool around the blockage
    \item \textbf{Reduced rectal capacity:} conditions that stiffen or reduce the rectum, including after surgery or radiotherapy
    \item \textbf{Neurological conditions:} stroke, multiple sclerosis, and spinal cord injury
    \item \textbf{Cognitive impairment:} dementia and other conditions
    \item \textbf{Other causes:} haemorrhoids, prolapse, and inflammatory bowel disease
\end{itemize}

\section*{Clinical Features}
The pattern of leakage gives clues to the cause.
\begin{itemize}
    \item Involuntary leakage of stool or wind
    \item Slight soiling of underwear with mucus or stool
    \item Urgency, with difficulty reaching the toilet in time
    \item Leakage after a normal bowel motion
    \item The need to wear pads
    \item Avoidance of social activities and travel
    \item Associated symptoms: constipation, diarrhoea, or pelvic symptoms
\end{itemize}

\section*{Differential Diagnosis}
Other conditions can be confused with faecal incontinence.
\begin{itemize}
    \item Diarrhoea from infection, irritable bowel syndrome, or medicines
    \item Haemorrhoids and anal fissures, which cause soiling
    \item Rectal prolapse
    \item Anal fistula or infection
    \item The distinction matters because the treatments differ
\end{itemize}

\section*{When to Seek Medical Care}
Anyone with faecal incontinence should discuss it with a GP, who can assess the cause and refer to a continence service or specialist. It is a treatable condition, and help is available.

\textbf{Contact your local emergency services for:}
\begin{itemize}
    \item Sudden incontinence with severe abdominal pain, vomiting, or fever
    \item Signs of a serious bowel problem, including a swollen, tender abdomen
    \item Blood in the stool with other worrying symptoms
    \item Sudden loss of bowel control with weakness, which may indicate a spinal emergency or stroke
\end{itemize}

\section*{Diagnosis and Investigations}
Assessment identifies the cause and its severity.
\begin{itemize}
    \item \textbf{History:} the pattern of leakage, bowel habit, childbirth, and surgery
    \item \textbf{Bowel diary:} recording leakage and bowel motions
    \item \textbf{Examination:} examination of the back passage and pelvic floor
    \item \textbf{Ultrasound:} to assess the anal sphincters
    \item \textbf{Anorectal manometry:} measures the pressure and function of the muscles
    \item \textbf{Rectal examination and endoscopy:} to exclude other bowel conditions
\end{itemize}

\section*{Management}
Management is effective and stepped.
\begin{itemize}
    \item \textbf{Pelvic floor exercises:} strengthening the sphincter muscles
    \item \textbf{Bowel habit training:} establishing regular toilet times
    \item \textbf{Dietary management:} adjusting fibre, fluids, and diet to firm the stool
    \item \textbf{Medicines:} to control diarrhoea or constipation
    \item \textbf{Continence products:} pads and aids for comfort and confidence
    \item \textbf{Biofeedback:} retraining the muscles with specialist guidance
    \item \textbf{Surgery:} sphincter repair, sacral nerve stimulation, or other procedures for severe cases
\end{itemize}

\section*{Complications}
\begin{itemize}
    \item Social isolation and withdrawal
    \item Depression and reduced self-esteem
    \item Skin irritation and infections around the back passage
    \item Reduced work and leisure participation
    \item Falls in older people rushing to the toilet
    \item Admission to residential care in some cases
\end{itemize}

\section*{Prognosis and Outcomes}
Most people with faecal incontinence improve substantially with treatment. Pelvic floor exercises, bowel management, and medicines help the majority, and surgical options are available for severe cases. The key to a good outcome is seeking help early and using the full range of support available.

\section*{Prevention and Self-Care}
\begin{itemize}
    \item Do pelvic floor exercises, especially after childbirth
    \item Maintain regular, healthy bowel habits
    \item Drink enough fluid and eat a balanced diet with appropriate fibre
    \item Plan toilet access and use the toilet before leaving home
    \item Use continence products without embarrassment
    \item Seek help early rather than accepting the problem
    \item Ask a doctor about referral to a continence service
\end{itemize}

\section*{Sources and Further Reading}
The following reputable sources provide further information for healthcare students and patients seeking reliable, current advice about faecal incontinence.
\begin{itemize}
    \item Continence Foundation of Australia. \textit{Bowel control and faecal incontinence.} https://www.continence.org.au/
    \item Healthdirect Australia. \textit{Bowel incontinence.} https://www.healthdirect.gov.au/bowel-incontinence
    \item NHS. \textit{Bowel incontinence.} https://www.nhs.uk/conditions/bowel-incontinence/
    \item Mayo Clinic. \textit{Fecal incontinence.} https://www.mayoclinic.org/diseases-conditions/fecal-incontinence/
    \item Jean Hailes for Women's Health. \textit{Pelvic floor health.} https://www.jeanhailes.org.au/
    \item MSD Manual. \textit{Fecal incontinence.} https://www.msdmanuals.com/
\end{itemize}
